%% =====================================================================
%%  CHAPTER 9: SUMMARY                                 (budget 2 pages)
%%  chapters/ch9-summary.tex
%% =====================================================================
\pcChapter{summary}{2}{סיכום}
\begin{pcguide}[מה כותבים כאן]
פרק הסיכום מחבר את הדוח לשורה תחתונה אחת: האם הפרויקט היתכן, על מה זה נשען, ומה עושים מיד לאחר האישור. הדוח הוא חוזה מול המרצה:
מה תעשו, מה תהיה התוצאה, ואיך ייבדק שעמדתם בכך. בדקו שאין סתירות בין הפרקים.
\end{pcguide}
\section{סינתזה: ארבעת עמודי ההיתכנות}\label{sec:synthesis}
הדוח בוחן אם כדאי לבנות אב-טיפוס מקומי של \latin{Grounded RAG} להפקת יחידות הדרכה ותרגול \threeD{} מחומרים פנימיים. ההיתכנות נשענת על ארבעה
עמודים, וטבלה~\ref{tab:pillars} מסכמת לכל אחד מהם את המסקנה, את ההוכחה בדוח ואת התנאי שנותר פתוח. הציונים והיעדים בדוח הם יעדי תכנון
ומודל חישובי; הם יאומתו במדידה בפיילוט, ואינם תוצאות שנצפו בתעשייה האווירית.
\begin{table}[!htbp]
\centering
\caption{ארבעת עמודי ההיתכנות: מסקנה, הוכחה בדוח ותנאי פתוח}\label{tab:pillars}
\begin{pctable}{H{0.17},H{0.37},H{0.22},H{0.24}}
  \pchead{עמוד היתכנות & מסקנה מסכמת & הוכחה בדוח & תנאי או שער פתוח}
  \pcrow{\textbf{צורך} \latin{(Need)} & ידע מקצועי ופעולות חפיפה תלויים במומחים בודדים; בארגון בהיקף \latin{15,000} עובדים בארבע חטיבות \cite{iai2025about} גם שיפור מתון בהכנת חומרי הדרכה רחב השפעה. הצורך נשען על הספרות ועל אפיון הבעיה ולא על נתוני \latin{IAI} & פרקים~\ref{ch:intro} ו-\ref{ch:problem}; \cite{zuin2025tacit} & מדידת בסיס ידני \latin{(D4)}; ראיונות מומחים (\risk{07})}
  \pcrow{\textbf{תוקף מדעי} \latin{(Scientific Validity)} & ספרות \latin{2023+} תומכת בכל מרכיב לחוד (אחזור מעוגן, הערכה פדגוגית אוטומטית, הדרכת \latin{XR}); הפער שנותר הוא השילוב: עיגון, אישור מומחה, תרחיש \threeDadj{} ומדידת זמן מול בסיס ידני & פרק~\ref{ch:lit}, סעיף~\ref{sec:gap}; \cite{hua2025civask} & בדיקת \hyp{1}--\hyp{4} בפיילוט; היעדים הם יעדי תכנון}
  \pcrow{\textbf{כדאיות טכנו-כלכלית} \latin{(Techno-Economic Viability)} & דגם מקומי בגודל \latin{7B}--\latin{8B} מכומת דורש בחישוב כ-\mth{9.4} \latin{GB} \latin{VRAM} (\mth{\approx12.6} \latin{GB} בהקשר ארוך), ולכן מתאים לכרטיס של \latin{24} \latin{GB}; שוק ה-\latin{SAM} כ-\money{\$}{4.62B} ו-\latin{SOM} של \money{\$}{76.5K}--\money{\$}{306K} לשנה: שוק צר אך מוגדר & פרקים~\ref{ch:market} ו-\ref{ch:alt}; טבלאות~\ref{tab:som} ו-\ref{tab:budget} & מדידת קצב יצירה בפועל (יעד \gesym{37} טוקנים לשנייה); זמינות \latin{GPU}}
  \pcrow{\textbf{החלטת ארכיטקטורה} \latin{(Architecture Decision)} & נבחרה חלופה ג (\num{3.95}) לעומת הבסיס הידני (\num{3.85}) וחלופת ה-\latin{LMS} (\num{3.35}). ההפרש קטן ורגיש למשקולות, ולכן הבחירה מותנית בשערי \latin{Go/No-Go} & טבלאות~\ref{tab:pugh} ו-\ref{tab:sens}; סעיף~\ref{sec:arch} & נאמנות \gesym{\pct{90}}, אפס בקשות לרשת חוץ, עמידה ב-\req{08}; בכשל: גיבוי ללא \latin{3D}}
\end{pctable}
\end{table}
המשמעות המנהלית: הבעיה אמיתית וניתנת לכימות, הספרות מצדיקה בדיקה, החומרה והעלות מצויות בהישג, והחלופה הנבחרת נבדקת בשערים שנקבעו מראש. מה שאינו
מוכח הוא הערך בתנאי הארגון: משך חפיפה, שיעור עזיבה והיקף הידע האבוד טרם נמסרו, ולכן יעדי הפיילוט (דיוק חילוץ \pctto{60}{90}, חיסכון בזמן
\gesym{\pct{50}} ונאמנות \gesym{\pct{90}}) מוצגים כיעדים שייבדקו.

\needspace{6\baselineskip}
\section{שלוש הפעולות המיידיות לאחר האישור}\label{sec:actions}
לאחר אישור המנחה נפתחים שלושה צעדים העוקבים זה אחר זה, והם מגדירים את תנאי הפתיחה של הפיילוט ב-\latin{M0} (איור~\ref{fig:actions}). הראשון הוא הקריטי
ביותר: בלי החלטה על גישה לנתונים אין קורפוס ארגוני, ולכן נקבע מראש מסלול חלופי על קורפוס פומבי (\risk{01}).
\begin{figure}[!htbp]
\centering
\pcFig{post_approval}{%
\pcAutoFit{\begin{pctikz}[x=1cm,y=1cm]
  \node[pcnodeD,text width=3.5cm,minimum height=1.0cm] (h1) at (0,0) {\tn{\latin{1.} בקשת נתונים והרשאות \latin{(Data/Access Request)}}};
  \node[pcnodeD,text width=3.5cm,minimum height=1.0cm] (h2) at (-4.3,0) {\tn{\latin{2.} הקמת \latin{GitHub} וכונן משותף \latin{(Setup)}}};
  \node[pcnodeD,text width=3.5cm,minimum height=1.0cm] (h3) at (-8.6,0) {\tn{\latin{3.} פגישת פתיחה \latin{(Kickoff Meeting)}}};
  \node[pcnode,text width=3.5cm,anchor=north,font=\scriptsize\sffamily] (d1) at (0,-0.8) {\tn{בקשה בכתב לאיש הקשר ב-\latin{IAI}: היקף הקורפוס, סיווג ומומחים לראיונות. אחראי: \StudentTwo; מיד לאחר האישור}};
  \node[pcnode,text width=3.5cm,anchor=north,font=\scriptsize\sffamily] (d2) at (-4.3,-0.8) {\tn{מאגר פרטי ותיקיות לקורפוס, שאלות אמת, יומנים ודוח; ללא מידע מסווג. אחראי: \StudentOne; עד \latin{M0}}};
  \node[pcnode,text width=3.5cm,anchor=north,font=\scriptsize\sffamily] (d3) at (-8.6,-0.8) {\tn{עם המנחה ואיש הקשר: אישור שאלת המחקר, הקורפוס והמדדים. אחראי: \StudentTwo; שבוע \latin{1} (\latin{07-03-2027})}};
  \draw[pcflow] (h1.west) -- (h2.east);
  \draw[pcflow] (h2.west) -- (h3.east);
  \node[pcnode,dashed,fill=white,text width=3.5cm,anchor=north,font=\scriptsize\sffamily] (fb) at (0,-3.6) {\tn{מסלול חלופי: אין אישור עד סוף שבוע \latin{2} $\Rightarrow$ קורפוס \latin{FAA} פומבי}};
  \draw[pcflow,dashed] (d1.south) -- (fb.north);
  \node[pcnode,fill=pcZebra,text width=3.5cm,anchor=north,font=\scriptsize\sffamily] (o3) at (-8.6,-3.6) {\tn{תוצר: היקף מאושר בכתב ופתיחת \latin{M0}}};
  \draw[pcflow] (d3.south) -- (o3.north);
\end{pctikz}}}
\caption{שלוש הפעולות המיידיות לאחר אישור הדוח: בקשת נתונים והרשאות, הקמת מאגר וכונן משותף ופגישת פתיחה; מועדים לפי לוח \latin{14} השבועות (פרק~\ref{ch:plan})}\label{fig:actions}
\end{figure}
\begin{pcbox}[pcGood]{החלטה מוצעת}
לאשר אב-טיפוס אקדמי מוגבל, ולא להתחייב לפריסה ארגונית לפני שיושגו מדדי הדיוק, הזמן, הבטיחות והאבטחה בפיילוט.
עמידה ביעדי הפיילוט אינה מהווה הסמכה מקצועית של עובד ואינה מהווה אישור אבטחה או בטיחות.
\end{pcbox}
\pcChapterEnd{summary}
